\section{Five Case Studies}
\label{sec:cases}

Each case study below is organized around three questions.
First, was the authority, on its statutory or constitutional face, intended as time-bounded?
Second, did the ratcheting pattern in fact occur,
in the sense that an authority granted as exceptional became routine or permanent?
Third, would the typed-rollback grammar developed in Part~\ref{sec:grammar} have structurally prevented the ratchet?
The third question is the most demanding and the one on which the Article's structural claim must stand or fall.
The Article is candid where the answer is contested or where the case study fits the framing less cleanly.
The five case studies vary in fit, and the variation is itself instructive.

\subsection{Weimar Article 48: the canonical case}
\label{sec:cases:weimar}

Article~48 of the Weimar Constitution \cite{TODO_weimar_const_art48} authorized the Reich President
to take ``the necessary measures'' to restore public security and order
in cases where it had been ``substantially disturbed or endangered.''
The article required the Reich President to notify the Reichstag of measures taken,
and permitted the Reichstag to demand the suspension of any such measure.
The drafters of the article, working within the constitutional convention at Weimar in 1919,
appear to have understood the authority as a narrowly bounded emergency power
designed to provide the new republic with a tool for response to acute civil disorder
\cite{TODO_caldwell_popular_sovereignty,TODO_jacobson_weimar_pdf}.

The historical trajectory of the article's use is well-documented.
Through the 1920s, the authority was invoked in response to discrete crises:
the hyperinflation of 1923, the political assassinations of the early decade,
and the recurring municipal-level civil disturbances that characterized the era.
The invocations during the early period were generally followed by Reichstag review,
and the measures taken were, in many cases, subsequently regularized through ordinary legislation
or allowed to lapse \cite{TODO_kershaw_hitler}.

The pattern shifted in the early 1930s.
The Bruening, Papen, and Schleicher governments of 1930 to 1933 used Article~48 not as a response to
discrete crises but as a substitute for ordinary legislation that could not pass the fragmented Reichstag.
By 1932, governance by Article~48 decree had become routine,
and the article's text -- which spoke of measures necessary to restore public order --
had been stretched to cover budgetary decisions, currency policy, and the regulation of commerce
\cite{TODO_caldwell_popular_sovereignty,TODO_mommsen_weimar}.
The Reichstag Fire Decree of February 28, 1933, issued under Article~48 in response to the burning of the Reichstag building,
suspended civil liberties indefinitely.
The Enabling Act of March 24, 1933, although procedurally distinct, was made possible in part by the
prior normalization of governance by emergency decree.

The historiography on whether the Article~48 ratchet was structural or political is contested.
Ian Kershaw \cite{TODO_kershaw_hitler} treats the trajectory as substantially political,
emphasizing the role of specific political actors and the contingency of particular decisions.
Hans Mommsen \cite{TODO_mommsen_weimar} treats it as more structural,
emphasizing the way that the article's text and the constitutional design surrounding it
provided no built-in mechanism for the lapse of decreed measures.
Peter Caldwell's work on popular sovereignty in Weimar
\cite{TODO_caldwell_popular_sovereignty} occupies an intermediate position.
The Article does not need to resolve this dispute.
It needs only the narrower observation that Article~48 contained no typed rollback obligation:
the Reich President could exercise the authority without exhibiting, at the moment of exercise,
a constructible path back to the prior state.
The decrees taken under the authority created facts on the ground -- suspended civil liberties,
imprisoned dissidents, dissolved political organizations -- that did not return to the prior state
when the formal emergency ended.

The typed-rollback grammar would have structurally constrained the article's use in a way the original
text did not.
A grammar that required the Reich President, at the moment of exercising Article~48 authority,
to exhibit a witness for the path back to the prior state would have made decrees creating irreversible
facts on the ground unconstructible.
The grammar would not have prevented the political use of the article to circumvent the Reichstag,
but it would have prevented the conversion of that political use into permanent constitutional facts.
This is the cleanest application of the framing to a historical case,
and the case study against which the Article's structural claim is most defensible.

\subsection{Section 230 and the takedown ecosystem}
\label{sec:cases:section230}

Section~230 of the Communications Decency Act \cite{TODO_section230} is not in its primary form an emergency-authority statute.
The provision immunizes interactive computer services from civil liability for content posted by users,
and immunizes good-faith content-moderation actions from civil liability.
The case-study selection here is consequently the one most in tension with the Article's framing.
The argument that Section~230 bears structural similarity to a delegated emergency authority requires acknowledgement
that the statute does not, on its face, grant any emergency power and does not, on its face, constitute a delegation
of executive authority.

The structural similarity is downstream of how the statute operates in conjunction with the Digital Millennium Copyright Act's
notice-and-takedown procedure \cite{TODO_dmca_512} and with the platform self-regulatory practices that have arisen
in the statute's shadow.
The DMCA notice-and-takedown procedure permits a copyright holder to request removal of allegedly infringing content.
The platform that receives the notice has a strong statutory incentive to remove the content without making an
independent merits determination, because removal preserves the safe harbor and refusal does not.
The platform that wishes to contest the notice has a counter-notice procedure available,
but the structural default of the regime is removal followed by the absence of any obligation to restore.
Section~230 immunizes the platform's good-faith removal decision regardless of whether the removal was substantively warranted.
The combined effect is a takedown regime in which removal is structurally a time-bounded executive act
of a private intermediary,
with no statutory mechanism that requires the intermediary to construct a path back to the prior state.\footnote{%
Daphne Keller has argued at length that the DMCA-Section~230 interaction produces a takedown regime that operates
in important respects like a privatized administrative regime
\cite{TODO_keller_takedowns}.
The Article's structural claim is consistent with Keller's empirical account.}

The ratcheting pattern in this regime is less dramatic than in the Weimar case but is documented.
Empirical work on DMCA notice-and-takedown indicates that a substantial fraction of notices target content for which
the underlying infringement claim is dubious or absent \cite{TODO_urban_quilter_dmca}.
The removed content is, in practice, not restored even when the counter-notice procedure produces no response from
the original requester.
The platform's safe harbor incentive favors removal as the default,
and the counter-notice procedure is rarely used because of the friction it imposes on the user whose content was removed
\cite{TODO_seng_dmca}.

The typed-rollback grammar would, in this regime, have a more specific implication.
A removal action under the DMCA-Section~230 regime would, under the grammar,
require the platform to construct, at the moment of removal, a witness for restoration of the content
in the case that the underlying claim is later determined to be without merit.
The witness is not the counter-notice procedure as currently constituted,
which operates at the moment of objection rather than at the moment of removal.
The witness is a precondition for the removal action itself.
A removal that cannot be reversed by construction would fail to type-check.

The application of the grammar to this regime is contested on multiple fronts.
Platforms operating at the scale of contemporary social-media services would argue that a typed-rollback discipline
is impractical at scale because the volume of removals is too high to permit individualized rollback construction.
Eric Goldman's writing on Section~230 \cite{TODO_goldman_section230} would observe that the grammar imposes a
structural cost on intermediaries that the existing statute deliberately declines to impose.
Danielle Citron's writing on intermediary liability \cite{TODO_citron_intermediary_liability} would argue, from a
different direction, that the typed-rollback discipline does not address the harms produced by content that should
have been removed but was not.
The Article acknowledges these objections without resolving them.
The structural similarity to a delegated emergency authority is the load-bearing observation;
the policy question of whether the grammar should apply to this regime is a separate question on which the
Article does not take a position.

\subsection{GDPR Article 17: erasure without rollback}
\label{sec:cases:gdpr}

Article~17 of the General Data Protection Regulation \cite{TODO_gdpr_art17} grants data subjects a right to erasure of
their personal data in specified circumstances.
The article's enforcement has produced a body of erasure orders against search-engine providers and platform operators,
of which the \emph{Google Spain} decision \cite{TODO_google_spain} is the canonical example.
The erasure orders are time-bounded executive acts in the formal sense that they are issued by administrative bodies
operating under defined procedures,
and they are subject to judicial review on familiar administrative-law grounds.

The ratcheting pattern in this regime takes a different form from the Weimar case.
There is no observed trajectory in which the erasure authority has been used to dissolve the underlying constitutional order
or to substitute for ordinary legislation.
The pattern that this Article identifies is the structural irreversibility of individual erasure orders.
An erasure order requires the controller to delete or de-index the targeted data.
The deletion is, as a practical matter, irreversible:
even if the order is later set aside by judicial review, the data is not, in general, restored
\cite{TODO_kuner_gdpr_erasure}.
The order is, in form, a time-bounded executive act.
It is, in effect, a permanent alteration of the data ecosystem.

The Article's framing here is narrower than in the Weimar case.
The claim is not that the GDPR Article~17 regime has ratcheted into something other than what it was designed to be.
The claim is that individual erasure orders share the structural defect of an emergency authority granted without a
typed rollback witness:
the order can be admitted without exhibiting, at the moment of admission, a constructible path back to the prior state.
The witness is not currently a requirement of the regime.

Mantelero's writing on GDPR enforcement \cite{TODO_mantelero_gdpr_enforcement} and Kuner's analysis of the right to be forgotten
\cite{TODO_kuner_gdpr_erasure} both observe the structural irreversibility of individual erasure orders.
Neither author proposes a typed-rollback discipline as the corrective.
The Article's contribution to this literature is the identification of the structural similarity to the emergency-authority
pattern,
and the proposal that the typed-rollback grammar provides a more principled corrective than ad hoc restoration mechanisms.

The fit of the case study to the framing is intermediate.
The GDPR Article~17 regime is not on its face an emergency-authority regime,
and the data-protection scholars who have written on it have not, in general, framed it as such.
The structural similarity is at the level of the form of authority rather than the substance of the regime.
The Article includes the case study because the structural similarity is real and because the corrective the
grammar proposes is applicable to the regime even though the regime is not, in its self-understanding, an emergency-authority regime.

\subsection{FISA Section 702: emergency surveillance authorizations}
\label{sec:cases:fisa702}

Section~702 of the Foreign Intelligence Surveillance Act, codified at 50 U.S.C. \S~1881a \cite{TODO_fisa_702_statute},
provides for the acquisition of foreign-intelligence information from non-U.S. persons reasonably believed to be located
outside the United States.
The section includes an emergency authorization mechanism that permits the Attorney General,
in conjunction with the Director of National Intelligence,
to authorize a surveillance program in advance of judicial approval where the Attorney General determines that an emergency exists.

The Privacy and Civil Liberties Oversight Board's 2014 report on the Section~702 program
\cite{TODO_pclob_702_report}
and the declassified opinions of the Foreign Intelligence Surveillance Court
\cite{TODO_fisc_702_opinions}
document the operation of the emergency authorization mechanism in practice.
The mechanism is invoked rarely in absolute terms but operates at a boundary of judicial review that the original
statutory text did not contemplate.
The information collected under an emergency authorization remains in the holdings of the collecting agency
even if the authorization is later determined to have been improperly granted.

The structural defect, on the framing this Article develops, is the absence of a typed rollback witness for
information collection.
The mechanism authorizes the collection;
it does not require, at the moment of authorization, the construction of a path back to the pre-collection state.
The targets of the collection do not return to a pre-collection state when the emergency authorization is later
reviewed or withdrawn.
The collected information persists.

The Article's claim here is the most cautious of the five case studies.
The public record on Section~702 emergency authorizations is incomplete.
The declassified opinions and PCLOB reports provide an authoritative but partial picture of the program's operation.
A national-security-law scholar with classified-clearance access to the full record may have grounds to dispute the
characterization above.
The structural claim is offered as a working hypothesis on the basis of the available public record,
with explicit acknowledgement that the record is incomplete.
The typed-rollback grammar's applicability to surveillance authorities depends on whether information collection can,
in principle, be made reversible by construction.
The cryptographic literature on revocable encryption \cite{TODO_revocable_encryption_lit} suggests that some forms of
collection can be made constructively reversible.
Whether the construction is operationally feasible in the surveillance context is an empirical question on which
the Article does not take a position.

\subsection{The 2001 AUMF: a quarter-century of unbounded authorization}
\label{sec:cases:aumf}

The Authorization for Use of Military Force enacted by Congress on September 18, 2001 \cite{TODO_aumf_2001}
permits the President to use ``all necessary and appropriate force'' against ``those nations, organizations, or persons
he determines planned, authorized, committed, or aided'' the September 11 attacks.
The text contains no geographic limit, no sunset, and no requirement of periodic reauthorization.

The trajectory of the AUMF's use is documented at length by the Brennan Center
\cite{TODO_brennan_center_aumf} and by scholarly accounts in Chesney \cite{TODO_chesney_aumf}
and Jaffer \cite{TODO_jaffer_aumf}.
By 2024, the authorization had been cited as legal basis for armed action in at least eight named theaters,
including territories of states not in existence as targets of policy at the time of the authorization's enactment.
The Islamic State of Iraq and the Levant, which did not exist in any form in 2001, has been characterized as a
covered organization under the AUMF on the legal theory that it is a successor entity to al-Qaeda in Iraq.
The expansion of the authority through the doctrine of associated forces has been documented in Department of Justice
white papers, declassified Office of Legal Counsel opinions, and the periodic ``war powers'' reports submitted to Congress
under the War Powers Resolution.

Periodic legislative efforts to repeal and replace the 2001 AUMF, including proposals by Senators Kaine and Young
across multiple Congresses, have not produced enacted legislation.
The pattern is the equilibrium failure described in Part~\ref{sec:pattern}:
the default at each successive period is continuation,
and the political coalition required to enact a replacement is the same coalition required to enact a repeal.
That coalition has not assembled in the twenty-five years since the original enactment.

The typed-rollback grammar's application to the AUMF is the most politically charged of the five case studies.
The application is structural: the original AUMF was enacted without a typed rollback witness,
which is to say, without a requirement that the President exhibit, at the moment of authorizing the use of force,
a constructible path back to the pre-authorization state.
The grammar would have required the original authorization to specify the conditions under which the use of force
was to terminate and the state of affairs to which the post-termination order would return.
A military action that could not exhibit, at the moment of authorization, a constructible path back to the prior state
would have failed to type-check.

The political content of this case study is not the Article's subject.
The Article does not take a position on the merits of any particular use of force under the AUMF or on the
desirability of any particular successor authorization.
The structural claim is that the original AUMF's grant of unbounded authority is a clear instance of the
ratcheting pattern,
and that the typed-rollback grammar would have made the ratchet structurally unavailable.
A reader who disagrees with the structural claim on this case study can disagree with it without disturbing the
Article's claims about the other four cases.
The AUMF is included because it is the most prominent contemporary American example of the ratcheting pattern,
and because a structural account that omitted it would be incomplete.
